\documentclass[11pt,a4paper]{article}

% proof box fix anyhow









%% ---------------------------------------------------------------- packages
\usepackage[a4paper,margin=2.5cm,headheight=14pt]{geometry}
\usepackage[T1]{fontenc}
% \usepackage{mathptmx}
\usepackage{microtype}
\usepackage{amsmath,amssymb,amsthm,mathtools}
\usepackage{xcolor}
\usepackage{enumitem}
\usepackage{booktabs,array,longtable,tabularx}
\usepackage{fancyhdr}
\usepackage{lastpage}
\usepackage[hidelinks]{hyperref}

\newcommand{\studentname}{Pranjal Solanki}
\newcommand{\rollno}{2025386}

%% -------------------------------------------------- shortcut macros
\newcommand{\N}{\mathbb{N}}
\newcommand{\Z}{\mathbb{Z}}
\newcommand{\Q}{\mathbb{Q}}
\newcommand{\R}{\mathbb{R}}
\newcommand{\pair}[2]{\{#1,#2\}}
\newcommand{\xor}{\oplus}
\newcommand{\imp}{\Longrightarrow}
\renewcommand{\iff}{\Longleftrightarrow}
\newcommand{\cmark}{\checkmark}

%% -------------------------------------------------- header / footer
% name in the top-left corner, page number (x of y) in the bottom-right corner
\pagestyle{fancy}
\fancyhf{}
\lhead{\studentname\ (\rollno)}
\rhead{Discrete Mathematics, Assignment 1}
\rfoot{Page \thepage\ of \pageref{LastPage}}
\renewcommand{\headrulewidth}{0.4pt}
\renewcommand{\footrulewidth}{0pt}
\fancypagestyle{plain}{\fancyhf{}\lhead{\studentname\ (\rollno)}\rfoot{Page \thepage\ of \pageref{LastPage}}\renewcommand{\headrulewidth}{0.4pt}}

\setlist[itemize]{itemsep=2pt,topsep=3pt}
\setlist[enumerate]{itemsep=2pt,topsep=3pt}
\setlength{\parindent}{0pt}
\setlength{\parskip}{5pt}
\renewcommand{\arraystretch}{1.15}

%% -------------------------------------------------- plain environments
% (simple replacements for the old coloured boxes)
\newenvironment{claimbox}[1][]{\par\smallskip\begin{quote}\textbf{#1.}\ \itshape}{\end{quote}}
\newenvironment{answerbox}[1][Answer]{\par\smallskip\begin{quote}\textbf{#1.}\ }{\end{quote}}
\newenvironment{verdictbox}[2][]{\par\begin{quote}}{\end{quote}}
\newenvironment{notebox}{\par\begin{quote}\textit{Remark.}\ }{\end{quote}}

\theoremstyle{definition}
\newtheorem*{lemma*}{Lemma}
\newtheorem*{fact*}{Fact}
% \renewcommand{symbol}{$\blacksquare$}

\newcommand{\pts}[1]{\hfill{\small\textit{[#1]}}}
\newcommand{\goodmark}{\textbf{Correct.}}
\newcommand{\badmark}{\textbf{Incorrect.}}
\newcommand{\textcolorplain}[2]{#2}

% simple section headings
\makeatletter
% \renewcommand\section{\@startsection{section}{1}{\z@}{-3.2ex \@plus -1ex \@minus -.2ex}{1.5ex \@plus .2ex}{\normalfont\large\bfseries}}
% \renewcommand\subsection{\@startsection{subsection}{2}{\z@}{-2.5ex \@plus -1ex \@minus -.2ex}{1ex \@plus .2ex}{\normalfont\normalsize\bfseries}}
% \renewcommand\subsubsection{\@startsection{subsubsection}{3}{\z@}{-2ex \@plus -1ex \@minus -.2ex}{0.8ex \@plus .2ex}{\normalfont\normalsize\bfseries\itshape}}
\makeatother
\renewcommand{\thesection}{Problem~\arabic{section}}
\renewcommand{\thesubsection}{\arabic{section}.\arabic{subsection}}
\renewcommand{\thesubsubsection}{\arabic{section}.\arabic{subsection}.\arabic{subsubsection}}
\renewcommand{\qedsymbol}{}
\begin{document}

\title{\textbf{Discrete Mathematics (CSE121)}\\[2pt] Assignment 1}
\author{\studentname \\ Roll No.\ \rollno \\ IIIT Delhi}
\date{}
\maketitle
\thispagestyle{fancy}

\noindent\textbf{Declaration.} 
I state that I used an AI tool (Claude.ai) 
to help prepare and typeset these solutions \textit{i used claude to generate shortcuts for various long 
mathematical commands to write them down in 1 single command ; the page numbers you see on the bottom and the top bar you see is all with help of claude guiding me thru my latex journey, i also plan on using claude for adding some colors in my pdf and i've used claude to beautify my solutions like the Q under bracket the different tables the checkmarks and all minor minor things}, 
and I didnt discuss solutions but had a chat regarding latex with \textit{Rudra Bhati and Paramveer}.

\medskip
\noindent\textbf{Summary of answers:}
\begin{center}
\begin{tabular}{@{ }ll@{ }}
\toprule
Problem & Result\\
\midrule
1(1) & My own personal experience with the 8 proof techniques\\
1(2a) & Claim false for $n=0$ (if $0\in\N$); proof is circular.\\
1(2b) & Claim false; the induction step needs $2^{-1}$ at $k=0$.\\
1(2c) & Proof is in fact correct.\\
1(2d) & Claim true; proof is wrong (divisibility confused with equality).\\
2.1 & Conversion of Statements.\\
2.2 & (a) tautology; (b) tautology; (c) tautology for odd $n$, contingency for even $n$.\\
3.1 & Minimum number of weighings is $4$.\\
3.2 & Minimum number of weighings is $5$.\\
4.1 & Infinitely many primes $\equiv3\pmod4$\\
4.2 & $\sqrt2+\sqrt3\notin\Q$.\\
5 & Smallest $N=12$; $A_n+B_n\le(5/4)^n$ for $n\ge10$.\\
\bottomrule
\end{tabular}
\end{center}

\section{}

\subsection{My experience with various proof techniques in this course.}
We have focused on various proof techniques as mentioned by sir in his pdf (when sir gives us 
the review before classes) I used the review as a base to start this question.

\begin{itemize}
  \item \textbf{Modus Ponens:}
        % (\ $[P \land (P \implies Q)] \implies Q$ ) 
        If the ``if'' statement is true, and its starting condition is true, then the result must be true. This is the best way to understand this

  \item \textbf{Proof by Counterexample:}
        (\ $\exists x \in U, \neg P(x)$ )
        when we are sure a specific case fails then we can disprove the entire statement ($\forall x, P(x)$) by showing the one specific case where the statement is False.

  \item \textbf{Direct Proof (Integers and Properties):}
        
        We directly assume the given statement $P$ is True, then we try to validate it by using various axioms, definitions and mathematics to show the conclusion $Q$ is true.

  \item \textbf{Proof by Contraposition:}
        (\ prove $P \implies Q$ by showing $\neg Q \implies \neg P$ ) we assume the negation of the conclusion is true, if proved true we have our proof because technically if negation is true the opposite must also be true.

  \item \textbf{Proving Biconditionals:}
        (\ $P \iff Q$ ) we split into two sections: first we prove the forward direction ( $P \implies Q$ ) and then we prove the opposite direction ($Q \implies P$). If both True we have our proof as True otherwise False.

  \item \textbf{Proof by Cases:}
        (\ $n$ is even or $n$ is odd ) we split our entire bigger statement into different smaller statements, if we can prove the smaller statements we can end up proving the bigger statement holds true overall.

  \item \textbf{Proof by Contradiction:}
        assume the entire statement is false, then show this assumption logically leads to an impossible mathematical absurdity.

  \item \textbf{Proof by Mathematical Induction:}
        (\ verify base case $P(b)$ for a boundary value such as $n=0$; assume inductive hypothesis $P(k)$, prove inductive step $P(k+1)$ ) establishes the validity of a statement across all integers.
\end{itemize}


At the end, the most important thing that I've learnt yet in this course is that a proof can be broken or built just from \emph{one single good decision}. 
To elaborate, you need to be sure about the method and if the method is correct the proof can be easily figured out just need to stay on the same path, 
its like a puzzle the more you look into it the more it either makes sense or makes your brain twist.

\subsection{Checking proposed proofs}


% here right now






\subsubsection*{(a)}
\begin{claimbox}[Claim]
If $n$ is a natural number and $n\mid n^2$, then $n\neq0$.
\end{claimbox}

\begin{verdictbox}{badred}
\badmark{} (The claim is True because obviously $0\notin\N$; The proof is invalid under every convention.)
\end{verdictbox}

\textbf{Incorrect part}
\begin{enumerate}
  \item The step \emph{``on cancelling $n$ from both sides, $n=k$''} is valid only if $n\neq0$ (cancellation
        $nx=ny\Rightarrow x=y$ needs $n\ne0$). But $n\neq0$ is exactly what we are trying to prove, so the this is absolutely incorrect to state.
  \item Even if we ignore that part and focus on $n=k$, the conclusion \emph{``therefore $n\neq 0$''} does not make any sense because how can we comment using that on $n$ being non-zero.
  \item In fact, the hypothesis $n\mid n^2$ is satisfied by \emph{every} $n$ (matter of fact $k=n$, since $n^2=n\cdot n$), so it can never force $n\ne0$.
\end{enumerate}

\textbf{Counterexample.} Take $n=0$. Then $n^2=0$ and $0=k\cdot 0$ for any integer $k$ (e.g.\ $k=5$), so $0\mid 0$ holds, yet $n=0$.
Hence the claim is false.
If instead $\N=\{1,2,3,\dots\}$, the conclusion $n\neq0$ is technically true because $n\ge1$ but that is only because of the definition of Natural Numbers $\N$, not from the hypothesis, and the given proof is still absolutely incorrect by 
mathematical standpoint.\\

\textbf{Fix.}
\begin{itemize}
  \item \emph{Corrected 1:} For every integer $n$, $n\mid n^2$. \emph{Proof:} $n^2=n\cdot n$, so $k=n$.
  \item \emph{Corrected 2:} If $n\neq0$ and $n^2=kn$ for an integer $k$, then $k=n$.
        \emph{Proof:} Since $n\neq0$ we may cancel $n$ (multiply by $n^{-1}$ in $\Q$): $n=k$.
\end{itemize}







\subsubsection*{(b)}
\begin{claimbox}[Claim]
$2^n=1$ for all non-negative integers $n$.
\end{claimbox}

\begin{verdictbox}{badred}
\badmark{} The claim is false (e.g.\ $2^1=2\ne1$); the inductive step is flawed.
\end{verdictbox}

\textbf{Incorrect part}
\[
2^{k+1}=\frac{2^k\cdot 2^k}{2^{k-1}} .
\]
This equation is true its equal $LHS=RHS$, but the proof then uses inductive hypothesis for $2^{k-1}$.
The hypothesis only covers exponents $m$ with $0\le m\le k$.
For $k\ge1$ this is fine, but for $k=0$ (i.e.\ proving $n=1$) the exponent $k-1=-1$ is \emph{outside} the range.
Because the step is invalid for $k=0$, the induction needs a \emph{proper new base case} $n=1$; that base case $2^1=1$ is false, so the induction never starts.
(Checking: the proof would give $2^1=\frac{1\cdot1}{2^{-1}}=\frac{1}{1/2}=2$, not $1$.)

\textbf{Fix.} Induction step that refers to $P(k-1)$ needs base cases $n=0$ and $n=1$ to be verified directly.

\begin{claimbox}[Corrected claim]
$2^n\ge1$ for all integers $n\ge0$, with equality if and only if $n=0$.
\begin{proof}
Induction on $n$.\\ \emph{Base Case} $n=0$: $2^0=1\ge1$.\\ \emph{Step Case:} 
if $2^k\ge1$ then $2^{k+1}=2\cdot2^k\ge2\ge1$, and in fact $2^{k+1}\ge2>1$, 
so equality occurs only at $n=0$.
\end{proof}

\end{claimbox}








\subsubsection*{(c)}
\begin{claimbox}[Claim]
Every positive integer $n$ can be written as a sum of distinct powers of $2$.
\end{claimbox}

\begin{verdictbox}{okgreen}
\goodmark{} This proof is \emph{valid}.
\end{verdictbox}

% hogya 1


\textbf{Confirmation.}
\begin{enumerate}
  \item \emph{Existence of $N$.} If the claim fails, the set of positive integers that are not sums of distinct 
  powers of $2$ is non-empty, so by the \textbf{well-ordering principle} 
  it has a least element $N$. \cmark
  \item \emph{$N>1$.} $1=2^0$ is a (one-term) sum of distinct powers of $2$, 
  so $N\neq1$; hence $N\ge2$. \cmark
  \item \emph{$N$ even, $N=2m$.} Then $m\ge1$ (as $N\ge2$) and $m<N$, so $m$ is 
  representable by minimality:
        $m=2^{a_1}+\dots+2^{a_r}$ with $a_1,\dots,a_r$ distinct. 
        Doubling, $N=2^{a_1+1}+\dots+2^{a_r+1}$, and the exponents
        $a_i+1$ are still distinct. So $N$ \emph{is} representable, 
        contradicting the choice of $N$. \cmark
  \item \emph{$N$ odd.} Since $N\ge2$ we get $N-1\ge1$, and $N-1<N$, so $N-1=2^{a_1}+\dots+2^{a_r}$ with distinct $a_i$ (minimality).
        \emph{Why is $a_i\ne0$ for all $i$?} (the one line the proof skips): $N-1$ is even. If some $a_i=0$, the sum would be
        $1+(\text{sum of terms }2^{a_j},\,a_j\ge1)$ = odd + even $=$ odd, a contradiction. So no term equals $2^0$.
        Therefore $N=2^0+2^{a_1}+\dots+2^{a_r}$ has pairwise distinct exponents $0,a_1,\dots,a_r$; $N$ is representable, a contradiction. \cmark
  \item Both cases contradict the choice of $N$, so no such $N$ exists. \cmark
\end{enumerate}

\textbf{Conclusion.} The proof is correct as proved by the above Confirmation.









\subsubsection*{(d)}
\begin{claimbox}[Claim]
Let $n$ be an integer. If $n^2+n$ is divisible by $2$, then $n$ is even or $n+1$ is even.
\end{claimbox}

\begin{verdictbox}{badred}
\badmark{} The claim is true, but the proposed proof is wrong.
\end{verdictbox}

\textbf{Incorrect part}
\begin{enumerate}
  \item \emph{``Since $2$ is prime, it cannot be expressed as a product of two integers 
  other than $1$ and $2$. Therefore one of the factors must equal $2$.''}. Divisibility means
        $n(n+1)=2k$ for some integer $k$, \emph{not} $n(n+1)=2$.
  \item The conclusion \emph{``$n=2$ or $n+1=2$''} is false: $n=3$ gives $n(n+1)=12$, divisible by $2$, but neither $3$ nor $4$ equals $2$.
  \item Also, \emph{``$n=2$ or $n+1=2$ $\Rightarrow$ $n$ or $n+1$ even''} would only cover $n\in\{2,1\}$, so even the final inference proves a much weaker claim than required.
\end{enumerate}

\textbf{Fix.} By the division algorithm, $n=2q$ or $n=2q+1$ for some integer $q$.
If $n=2q$, then $n$ is even. If $n=2q+1$, then $n+1=2(q+1)$ is even. Hence for \emph{every} integer $n$ or $n+1$ is even
(so the hypothesis $2\mid n^2+n$ is automatically true). 






% ye bhosdika deekhna hai


%% =====================================================================
%%                               PROBLEM 2
%% =====================================================================
\section{}

\subsection{Conversion of statements}

\subsubsection{}
\textbf{Setup.} Let $G=(V,E)$ be a graph, $\mathrm{Adj}(x,y)$ mean ``$x,y$ share an edge'', and let $\mathrm{CC}(G)$ denote ``$G$ is completely connected'' 
As provided graph is completely connected if and only if for every pair of distinct vertices, there exists a third
vertex that shares an edge with both, unless they already share a direct edge.

\textbf{Logical Expression.}
\[
\Phi\;:\equiv\;\mathrm{CC}(G)\;\Longleftrightarrow\;\underbrace{\forall x\,\forall y\,\Big[\big(x\neq y\wedge\neg\mathrm{Adj}(x,y)\big)\longrightarrow\exists z\,\big(z\neq x\wedge z\neq y\wedge\mathrm{Adj}(x,z)\wedge\mathrm{Adj}(y,z)\big)\Big]}_{Q}
\]

\textbf{Negation.} $P:=\mathrm{CC}(G)$, so $\Phi=P\iff Q$.
\begin{align*}
\neg(P\iff Q)&\equiv(P\wedge\neg Q)\vee(\neg P\wedge Q)\\[-2pt]
&\qquad\small \text{(a biconditional fails if and only if exactly one side holds)}\\[2pt]
\neg Q&\equiv\exists x\,\exists y\,\neg\Big[(x\neq y\wedge\neg\mathrm{Adj}(x,y))\to\exists z(\dots)\Big]\\[-2pt]
&\qquad\small (\neg\forall\equiv\exists\neg)\\
&\equiv\exists x\,\exists y\,\Big[(x\neq y\wedge\neg\mathrm{Adj}(x,y))\wedge\neg\exists z(\dots)\Big]\\[-2pt]
&\qquad\small (\neg(A\to B)\equiv A\wedge\neg B)\\
&\equiv\exists x\,\exists y\,\Big[x\neq y\wedge\neg\mathrm{Adj}(x,y)\wedge\forall z\,\big((z\neq x\wedge z\neq y)\to(\neg\mathrm{Adj}(x,z)\vee\neg\mathrm{Adj}(y,z))\big)\Big]\\[-2pt]
&\qquad\small (\neg\exists\equiv\forall\neg;\ \text{De Morgan})
\end{align*}
Hence
\[
\neg\Phi\equiv\Big(\mathrm{CC}(G)\wedge\neg Q\Big)\ \vee\ \Big(\neg\mathrm{CC}(G)\wedge Q\Big),
\]
with $\neg Q$ as displayed above.

\textbf{In English.} \emph{Either} $G$ is completely connected, yet there exist two distinct, non-adjacent vertices $x,y$ such that every other vertex fails to be adjacent to at least one of them (no common neighbour);
\emph{or} $G$ is not completely connected, yet every two distinct non-adjacent vertices have a third vertex adjacent to both.







% ye bhosdika deekhna hai ye


\subsubsection{}
\textbf{Definitions.} $\mathrm{Book}(b)$, $\mathrm{Reader}(r)$, $\mathrm{Chapter}(c)$, $\mathrm{Auth}(a,b)$ ($a$ is the author of book $b$), $\mathrm{Wrote}(a,c)$ ($a$ wrote chapter $c$), $\mathrm{Appr}(r,c)$ ($r$ appreciates $c$).

\textbf{Logical Expression.}
\[
\forall b\,\Big[\mathrm{Book}(b)\to\exists r\,\Big(\mathrm{Reader}(r)\wedge\forall a\,\forall c\,\big((\mathrm{Auth}(a,b)\wedge\mathrm{Chapter}(c)\wedge\mathrm{Wrote}(a,c))\to\mathrm{Appr}(r,c)\big)\Big)\Big]
\]

\textbf{Negation.}
\begin{align*}
&\neg\forall b\,[\mathrm{Book}(b)\to\exists r(\dots)]\\
&\equiv\exists b\,[\mathrm{Book}(b)\wedge\neg\exists r(\mathrm{Reader}(r)\wedge\forall a\forall c(\dots))]\\[-2pt]
&\qquad\small (\neg\forall,\ \neg(A\to B)\equiv A\wedge\neg B)\\
&\equiv\exists b\,[\mathrm{Book}(b)\wedge\forall r(\mathrm{Reader}(r)\to\neg\forall a\forall c(\dots))]\\[-2pt]
&\qquad\small (\neg\exists\equiv\forall\neg;\ \neg(A\wedge B)\equiv A\to\neg B)\\
&\equiv\exists b\,\Big[\mathrm{Book}(b)\wedge\forall r\,\Big(\mathrm{Reader}(r)\to\exists a\,\exists c\,\big(\mathrm{Auth}(a,b)\wedge\mathrm{Chapter}(c)\wedge\mathrm{Wrote}(a,c)\wedge\neg\mathrm{Appr}(r,c)\big)\Big)\Big]
\end{align*}
\textbf{In English.} There is a book such that \emph{every} reader fails to appreciate at least one chapter written by that book's author.





% padhliya ye 






\subsubsection{}
\textbf{Definitions.} $\mathrm{St}(s)$, $\mathrm{Took}(s,c)$, $\mathrm{A}(s,c)$ ($s$ received an A in class $c$). Let $G(s):\equiv\forall c\,(\mathrm{Took}(s,c)\to\mathrm{A}(s,c))$ (``all-A student'').

\textbf{Logical Expression.}
\[
\exists s\,\Big[\mathrm{St}(s)\wedge G(s)\wedge\forall t\,\big((\mathrm{St}(t)\wedge G(t))\to t=s\big)\Big]
\]
\textbf{Negation.}
\begin{align*}
&\forall s\,\Big[\neg\mathrm{St}(s)\vee\neg G(s)\vee\neg\forall t\big((\mathrm{St}(t)\wedge G(t))\to t=s\big)\Big]\\[-2pt]
&\qquad\small (\neg\exists\equiv\forall\neg;\ \text{De Morgan})\\
&\equiv\forall s\,\Big[\mathrm{St}(s)\to\Big(\neg G(s)\ \vee\ \exists t\,\big(\mathrm{St}(t)\wedge G(t)\wedge t\neq s\big)\Big)\Big]\\
&\equiv\forall s\,\Big[\mathrm{St}(s)\to\Big(\exists c\,(\mathrm{Took}(s,c)\wedge\neg\mathrm{A}(s,c))\ \vee\ \exists t\,\big(\mathrm{St}(t)\wedge t\neq s\wedge\forall c\,(\mathrm{Took}(t,c)\to\mathrm{A}(t,c))\big)\Big)\Big]
\end{align*}
\textbf{In English.} Every student either received a non-A in some class they took, or there is a \emph{different} student who received an A in every class they took. % Equivalently: \emph{either no student got all As, or at least two did.}

% \begin{notebox}
% A student who took no classes satisfies $G(s)$ vacuously; if this is unwanted, add $\exists c\,\mathrm{Took}(s,c)$ to $G(s)$. The negation changes accordingly.
% \end{notebox}











\subsection{Tautology, contradiction or contingency}

\subsubsection{(a)}
\[
u:=\big((p_1\Rightarrow p_2)\wedge(p_2\Rightarrow p_3)\wedge\dots\wedge(p_{n-1}\Rightarrow p_n)\wedge(p_n\Rightarrow p_1)\big)\Rightarrow(p_1\iff p_n).
\]
\begin{answerbox}[Answer]
$u$ is a \textbf{tautology} for every $n\ge3$.
\end{answerbox}


\begin{proof} Below mentioned proves my statement
\begin{itemize}
  \item Suppose $p_1$ is true. Then $p_1\Rightarrow p_2$ forces $p_2$ true, then $p_2\Rightarrow p_3$ forces $p_3$ true, \dots, and finally $p_{n-1}\Rightarrow p_n$ forces $p_n$ true (a chain of modus ponens). So $p_1,p_n$ are both true and $p_1\iff p_n$ holds. (this was also given in midsem)
  \item Suppose $p_1$ is false. If $p_n$ were true, then $p_n\Rightarrow p_1$ would be false, contradicting $p_1\iff p_n$. So $p_n$ is false and $p_1\iff p_n$ holds (both false).
\end{itemize}
In either case the $u$ holds, so $u$ is true under all assignments.
\end{proof}








\subsubsection{(b)}
\[
u:=\big(p_1\Rightarrow(p_2\Rightarrow(p_3\Rightarrow\dots\Rightarrow p_n))\big)\Rightarrow\big((p_1\wedge\dots\wedge p_{n-1})\Rightarrow p_n\big).
\]
\begin{answerbox}[Answer]
$u$ is a \textbf{tautology}.
\end{answerbox}
Let $F_k:=p_k\Rightarrow(p_{k+1}\Rightarrow(\dots\Rightarrow p_n))$ for $1\le k\le n-1$, so $F_{n-1}=p_{n-1}\Rightarrow p_n$ and $F_k=p_k\Rightarrow F_{k+1}$.

\begin{lemma*}
$F_k$ is false $\iff$ $p_k,p_{k+1},\dots,p_{n-1}$ are all true and $p_n$ is false.
\end{lemma*}
\begin{proof}
Induction on $k$.\\ \emph{Base} $k=n-1$: $p_{n-1}\Rightarrow p_n$ is false if and only if $p_{n-1}$ true and $p_n$ false.\\ \emph{Step:} $F_k=p_k\Rightarrow F_{k+1}$ is false if and only if $p_k$ is true and $F_{k+1}$ is false, if and only if (induction hypothesis) $p_k$ true and $p_{k+1},\dots,p_{n-1}$ true, $p_n$ false.
\end{proof}
Now let $L:=F_1$ and $R:=(p_1\wedge\dots\wedge p_{n-1})\Rightarrow p_n$. $R$ is false if and only if $p_1,\dots,p_{n-1}$ are all true and $p_n$ is false, which by the Lemma (with $k=1$) is exactly when $L$ is false. So $L\equiv R$. Then $u=L\Rightarrow R$ is $L\Rightarrow L$, which is true for every assignment: hence, a tautology. 





% ye check karna hai bhosdika

\subsubsection{(c)}
\[
u:=\big((p_1\oplus p_2)\wedge(p_2\oplus p_3)\wedge\dots\wedge(p_n\oplus p_1)\big)\Rightarrow(p_1\wedge\neg p_1).
\]

\begin{answerbox}[Answer]
$u$ is also a tautology for odd $n$, contingency for even $n$.
\end{answerbox}
\textbf{Key observation.} $p_1\wedge\neg p_1$ is always false, so $A\Rightarrow\text{False}\equiv\neg A$. Thus
$u\equiv\neg\big[(p_1\oplus p_2)\wedge\dots\wedge(p_n\oplus p_1)\big]$, and $u$ is false exactly when all consecutive pairs differ: $p_1\ne p_2\ne p_3\ne\dots\ne p_n\ne p_1$.

\textbf{When can this happen?} If $p_i\neq p_{i+1}$ for all $i$, the truth values alternate as: $p_1=T,\,p_2=F,\,p_3=T,\dots$, i.e.\ $p_i=p_1$ iff $i$ is odd. 
The last condition $p_n\neq p_1$ then requires $p_n$ to differ from $p_1$, i.e.\ $n$ must be \emph{even}. Formally, going around the cycle flips the value $n$ times and must return to the start, which forces $n$ to be even.
\begin{itemize}
  \item \textbf{$n$ odd:} statement is never true, so $u$ is never false: \textbf{tautology}.
  \item \textbf{$n$ even:} the alternating assignment $(T,F,T,F,\dots,T,F)$ makes the statement true and hence $u$ false; the all-true assignment makes the statement false and $u$ true. So $u$ is a \textbf{contingency}.
\end{itemize}













%% =====================================================================
%%                               PROBLEM 3
%% =====================================================================
\section{}

% \paragraph{A counting lemma used in both parts.}
% A weighing has exactly three outcomes ($<,=,>$). Let a \emph{strategy} be any (possibly adaptive) procedure, where the pans used in a weighing may depend on earlier outcomes.

% \begin{claimbox}[Lemma (information bound)]
% A strategy using at most $k$ weighings can produce at most $3^k$ different sequences of outcomes. Hence if there are $M$ scenarios that must be told apart and $M>3^k$, no strategy with $k$ weighings exists.
% \end{claimbox}
% \begin{proof}
% Induction on $k$. For $k=0$ there is exactly one (empty) sequence $=3^0$. If $k\ge1$, the first weighing is fixed and has $3$ possible outcomes; after each, the rest is a strategy with at most $k-1$ weighings having at most $3^{k-1}$ sequences, giving $\le3\cdot3^{k-1}=3^k$ in total.
% For the second claim: fix concrete weights for each scenario (genuine coin weight $1$, counterfeit $1\pm\tfrac12$ as appropriate). Then the outcome sequence is a \emph{function} of the scenario. If $M>3^k$, by the pigeonhole principle two different scenarios give the same outcome sequence; the strategy would then output the same answer for both and be wrong for at least one.
% \end{proof}

%% ---------------------------------------------------------------- 3.1
\subsection{9 coins, 2 lighter counterfeits}

\begin{answerbox}
The minimum number of weighings is \textbf{4}. (A question similar to Question 3 was in midsem and taught in class but I couldnt still grasp sections so i used AI for asistance in this question as it seems like reading the slides didnt help me figure it out)
\end{answerbox}

\subsubsection*{Lower bound: 3 weighings are not enough}
A possibility is just the \emph{pair} of counterfeit coins, so there are $\binom92=36$ scenarios. With $k\le3$ weighings there are at most $3^3=27<36$ outcome sequences. By the Lemma, no strategy with $3$ (or fewer) weighings works. 



\subsubsection*{Upper bound: a strategy with 4 weighings}
Labeling the coins $1,\dots,9$ and put $A=\{1,2,3\}$, $B=\{4,5,6\}$, $C=\{7,8,9\}$. Write $L<R$ if the left pan is lighter. 
% Throughout, ``$\{i,j\}$'' means ``coins $i,j$ are the counterfeit ones''. 
The counterfeit coins are identical, so the number of counterfeit coins on a pan determines which pan is lighter (equal pan sizes are always used).

\textbf{Weighing 1: $A$ v/s $B$.} Let $a,b,c$ be the number of counterfeit coins in $A,B,C$ ($a+b+c=2$). The pans balance iff $a=b$.

\begin{enumerate}[label=\textbf{Case \arabic*.},leftmargin=*,itemsep=6pt]
\item \textbf{$A<B$ (so $a>b$).} Then $(a,b,c)\in\{(2,0,0),(1,0,1)\}$, giving $3+9=12$ candidates:
 $\{12,13,23\}\cup\{ac: a\in A,c\in C\}$. All coins of $B$ are genuine; use coin $4$ as a reference genuine coin.

\textbf{Weighing 2: $\{1,7\}$ vs $\{2,8\}$.}
\begin{center}\small
\renewcommand{\arraystretch}{1.2}
\begin{tabular}{@{}lll@{}}
\toprule
Outcome & Remaining candidates & Weighing 3, then 4 \\
\midrule
$=$ & $\{12,\,18,\,27,\,39\}$ & W3: $1$ vs $4$. $<$: $\{12,18\}$; $=$: $\{27,39\}$.\\
& & W4 (both groups): $2$ vs $4$. $<$: $\{12\}$ resp.\ $\{27\}$; $=$: $\{18\}$ resp.\ $\{39\}$.\\[2pt]
$<$ & $\{13,\,17,\,19,\,37\}$ & W3: $1$ vs $3$. $=$: $\{13\}$; $<$: $\{17,19\}$; $>$: $\{37\}$.\\
& & W4 (for $\{17,19\}$): $7$ vs $4$. $<$: $\{17\}$; $=$: $\{19\}$.\\[2pt]
$>$ & $\{23,\,28,\,29,\,38\}$ & W3: $2$ vs $3$. $=$: $\{23\}$; $<$: $\{28,29\}$; $>$: $\{38\}$.\\
& & W4 (for $\{28,29\}$): $8$ vs $4$. $<$: $\{28\}$; $=$: $\{29\}$.\\
\bottomrule
\end{tabular}
\end{center}
\textit{Why the splits are correct (spot checks).} In W2 the pair $\{12\}$ puts one counterfeit on each pan, so it balances; $\{13\}$ has coin $1$ on the left only, so left is lighter; $\{23\}$ has coin $2$ on the right only, so right is lighter; $\{17\}$ has two on the left $\Rightarrow$ left lighter; $\{28\}$ has two on the right $\Rightarrow$ right lighter; the other pairs are checked identically. In W3 ($=$ branch), $1$ vs $4$: $\{12\},\{18\}$ contain coin $1$ (left lighter) while $\{27\},\{39\}$ do not (balance). Every leaf has exactly one candidate, so $4$ weighings suffice.

\item \textbf{$B<A$.} Symmetric to Case~1 with the roles $A\leftrightarrow B$ ($1,2,3\leftrightarrow4,5,6$) exchanged and coin $1$ as the genuine reference coin.

\item \textbf{$A=B$ (so $a=b$).} Either $(a,b,c)=(0,0,2)$: $\{78,79,89\}$, or $(1,1,0)$: $\{ab\}$ with $a\in A,b\in B$ (9 pairs). Total $12$ candidates.

 \textbf{Weighing 2: $\{1,4,7\}$ vs $\{2,5,8\}$.}
 \begin{center}\small
 \renewcommand{\arraystretch}{1.2}
 \begin{tabular}{@{}lll@{}}
 \toprule
 Outcome & Remaining candidates & Weighing 3, then 4 \\
 \midrule
 $=$ & $\{78,\,15,\,24,\,36\}$ & W3: $1$ vs $2$. $<$: $\{15\}$; $>$: $\{24\}$; $=$: $\{78,36\}$.\\
 & & W4 (for $\{78,36\}$): $7$ vs $3$. $<$: $\{78\}$; $>$: $\{36\}$.\\[2pt]
 $<$ & $\{79,\,14,\,16,\,34\}$ & W3: $1$ vs $3$. $=$: $\{79\}$; $<$: $\{14,16\}$; $>$: $\{34\}$.\\
 & & W4 (for $\{14,16\}$): $4$ vs $6$. $<$: $\{14\}$; $>$: $\{16\}$.\\[2pt]
 $>$ & $\{89,\,25,\,26,\,35\}$ & W3: $2$ vs $3$. $=$: $\{89\}$; $<$: $\{25,26\}$; $>$: $\{35\}$.\\
 & & W4 (for $\{25,26\}$): $5$ vs $6$. $<$: $\{25\}$; $>$: $\{26\}$.\\
 \bottomrule
 \end{tabular}
 \end{center}
\end{enumerate}
In every branch the two counterfeit coins are identified after at most $4$ weighings (I also checked all $36$ pairs by running this strategy by computer). Together with the lower bound, the minimum is exactly $4$. 



% same as previous bhosdika


%% ---------------------------------------------------------------- 3.2
\subsection{Twelve coins, two equal counterfeits, heavier or lighter unknown}

\begin{answerbox}
The minimum number of weighings is \textbf{5}.
\end{answerbox}

\subsubsection*{Lower bound: 4 weighings are not enough}
A scenario consists of (i) the pair of counterfeit coins: $\binom{12}{2}=66$ choices, and (ii) whether they are both heavier or both lighter: $2$ choices. By the rule of product there are $66\cdot2=132$ scenarios, all of which must be distinguished (the answer must name the coins \emph{and} the direction).
With $4$ weighings there are at most $3^4=81<132$ outcome sequences, so by the Lemma $4$ weighings (or fewer) cannot suffice. 

\subsubsection*{Upper bound: a strategy with 5 weighings}
\textbf{Weighing 1: $\{1,2,4,11\}$ vs $\{5,6,9,10\}$} (any $4$ vs $4$ works by symmetry). Call the pans $P$ (left), $Q$ (right), and the four unused coins $E$. Let a pair have $(l,r,e)$ counterfeit coins on $(P,Q,E)$:

\begin{center}\small
\renewcommand{\arraystretch}{1.2}
\begin{tabular}{@{}ccccc@{}}
\toprule
$(l,r,e)$ & number of pairs & if counterfeit heavy & if counterfeit light & scenarios \\
\midrule
$(2,0,0)$ & $\binom42=6$ & $P>Q$ & $P<Q$ & 12\\
$(0,2,0)$ & $6$ & $P<Q$ & $P>Q$ & 12\\
$(1,1,0)$ & $4\cdot4=16$ & $=$ & $=$ & 32\\
$(1,0,1)$ & $4\cdot4=16$ & $P>Q$ & $P<Q$ & 32\\
$(0,1,1)$ & $16$ & $P<Q$ & $P>Q$ & 32\\
$(0,0,2)$ & $6$ & $=$ & $=$ & 12\\
\bottomrule
\end{tabular}
\end{center}
So the three outcomes leave
\[
P>Q:\ 6+6+16+16=44,\qquad P<Q:\ 44,\qquad P=Q:\ 32+12=44
\]
scenarios, each $\le 3^4=81$, so the counting bound permits a continuation with $4$ more weighings. Hence, $4$ weighings are impossible ($132>81$).





% ye bhosdika deeekh 



%% =====================================================================
%%                               PROBLEM 4
%% =====================================================================
\section{}

\subsection{Infinitely many primes of the form $4k+1$}

\begin{claimbox}[Claim]
There are infinitely many primes $p$ with $p\equiv1\pmod4$.
\end{claimbox}

\textbf{Euclid's Lemma (of $x^2+1$)} If an odd prime $q$ divides $x^2+1$ for some integer $x$, then $q$ leaves a remainder of $1$ when divided by $4$. Equivalently, no odd prime that leaves a remainder of $3$ when divided by $4$ can divide a number of the form $x^2+1$.

\begin{proof}
\emph{Suppose, for contradiction,} that there are only finitely many primes that leave a remainder of $1$ when divided by $4$. The list is non-empty because $5$ is one, so write it as $p_1=5,p_2,\dots,p_r$ ($r\ge1$). Let
\[
P:=p_1p_2\cdots p_r,\qquad N:=(2P)^2+1 = 4P^2+1 .
\]
\begin{enumerate}
  \item \emph{$N>1$ and $N$ is odd.} $P\ge5$ so $N\ge101$, and $N=4P^2+1$ is odd.
  \item \emph{$N$ leaves a remainder of $1$ when divided by $4$.} Indeed, $N=4(P^2)+1$.
  \item \emph{$N$ has a prime factor that leaves a remainder of $1$ when divided by $4$.} $N>1$ so it factors into primes $q_1\cdots q_s$. Since $N$ is odd, all its prime factors $q_j$ must be odd. By construction, each $q_j$ divides $N = (2P)^2+1$. According to the lemma, any odd prime dividing a number of the form $x^2+1$ must leave a remainder of $1$ when divided by $4$. Therefore, every prime factor $q$ of $N$ leaves a remainder of $1$ when divided by $4$.
  \item \emph{$q$ is on the list.} By our assumption the list contains \emph{all} primes leaving a remainder of $1$ when divided by $4$, so $q=p_i$ for some $i$, hence $q\mid P$ and $q\mid 2P$, which means $q\mid (2P)^2$.
  \item \emph{Contradiction.} Now $q\mid N=(2P)^2+1$ and $q\mid(2P)^2$, so $q\mid N-(2P)^2=1$, impossible for a prime.
\end{enumerate}
The assumption is false, so there are infinitely many primes that leave a remainder of $1$ when divided by $4$. 
\end{proof}














\subsection{$\sqrt2+\sqrt3$ is irrational}

\begin{claimbox}[Lemma ($\sqrt2$ is irrational)]
$\sqrt2\notin\Q$.
\end{claimbox}
\begin{proof}
Suppose $\sqrt2=a/b$ with integers $a,b$, $b\neq0$, such that ($\gcd(a,b)=1$). Then $a^2=2b^2$, so $a^2$ is even. If $a$ were odd, $a=2m+1$ gives $a^2=4m^2+4m+1$ odd; hence $a$ is even, $a=2c$. Then $4c^2=2b^2$, i.e.\ $b^2=2c^2$, so $b$ is even by the same argument. Thus $2\mid\gcd(a,b)$, contradicting $\gcd(a,b)=1$.
\end{proof}

\begin{claimbox}[Claim]
$\sqrt2+\sqrt3$ is irrational.
\end{claimbox}
\begin{proof}
Suppose, for contradiction, that \[ \sqrt2+\sqrt3=\frac pq, \] where $p,q\in\mathbb Z$, $q\neq0$, and $\gcd(p,q)=1$. Since $\sqrt2+\sqrt3>0$, we have $p\neq0$.\begin{align*}
\sqrt3&=\frac pq-\sqrt2 &&\text{(rearranging)}\\ 3&=\frac{p^2}{q^2}-\frac{2p}{q}\sqrt2+2 &&\text{(squaring both sides)}\\ 1&=\frac{p^2}{q^2}-\frac{2p}{q}\sqrt2 &&\text{(rearranging)}\\ \frac{2p}{q}\sqrt2&=\frac{p^2-q^2}{q^2} &&\text{(rearranging)}\\ \sqrt2&=\frac{p^2-q^2}{2pq} &&\text{(dividing by }2p/q\neq0\text{)}.
\end{align*}
Since $p,q\in\mathbb Z$ and $p,q\neq0$, the right-hand side is rational. Hence $\sqrt2\in\mathbb Q$, contradicting the Lemma. Therefore, $\sqrt2+\sqrt3$ is irrational.\end{proof}









%% =====================================================================
%%                               PROBLEM 5
%% =====================================================================
% ---------- Problem 5, simple wording ----------
% Paste this as your 5th section.

\section{}

% \textbf{Setting up.}
Robot $A$ jumps $2$ or $5$ steps, and robot $B$ jumps $3$ or $7$ steps. The order of jumps matters
(for example $2$ then $5$ is a different way from $5$ then $2$).
Let $A_0=B_0=1$ (standing at step $0$ is one way: do nothing) and let $A_n=B_n=0$ for $n<0$. (I've also used AI to figure out the best way to present and writeout this question because its similar to midsem and classroom slide question and i am sure i was on the right path i just couldnt properly present it)

Thinking about the \emph{last} jump. To land on step $n$, robot $A$ either jumped $2$ from step $n-2$, or jumped $5$ from step $n-5$. So 
\[
A_n=A_{n-2}+A_{n-5},\qquad B_n=B_{n-3}+B_{n-7}\qquad(n\ge1).
\]
Using these, I wrote the initial values in a table (we need them later):
\begin{center}\small
\renewcommand{\arraystretch}{1.15}
\setlength{\tabcolsep}{4.5pt}
\begin{tabular}{@{}l|*{21}{r}@{}}
$n$ & 0&1&2&3&4&5&6&7&8&9&10&11&12&13&14&15&16&17&18&19&20\\\hline
$A_n$ & 1&0&1&0&1&1&1&2&1&3&2&4&4&5&7&7&11&11&16&18&22\\
$B_n$ & 1&0&0&1&0&0&1&1&0&1&2&0&1&3&1&1&4&3&1&5&5\\
\end{tabular}
\end{center}
% ``A robot can reach step $n$'' just means its count is at least $1$.

\subsection{Is there an $N$ after which both robots can reach every step?}
\begin{answerbox}
Yes. Take $N=12$: every step $n\ge12$ can be reached by both robots. Listed out in the table/
\end{answerbox}

\textbf{Robot $A$ (jumps 2 and 5) can reach every $n\ge4$.}
\begin{itemize}
  \item If $n$ is even and $n\ge4$: do $n/2$ jumps of size $2$.
  \item If $n$ is odd and $n\ge5$: do one jump of $5$, then the rest is $n-5$, which is even and $\ge0$, so do $(n-5)/2$ jumps of size $2$.
\end{itemize}

\textbf{Robot $B$ (jumps 3 and 7) can reach every $n\ge12$.} Look at the remainder of $n$ when divided by $3$.
\begin{itemize}
  \item $n=3q$: just do $q$ jumps of size $3$.
  \item $n=3q+1$: do one jump of $7$ and then $q-2$ jumps of $3$, because $7+3(q-2)=3q+1$. This needs $q\ge2$, i.e.\ $n\ge7$, which is true for $n\ge12$.
  \item $n=3q+2$: do two jumps of $7$ and then $q-4$ jumps of $3$, because $14+3(q-4)=3q+2$. This needs $q\ge4$, i.e.\ $n\ge14$.
        If $n\ge12$ and $n$ leaves remainder $2$ when divided by $3$, then the smallest such $n$ is $14$, so we are fine.
\end{itemize}

Since $12\ge4$, both robots can reach every $n\ge12$. \qed

\subsection{The smallest such $N$}
\begin{answerbox}
The smallest value is $N=12$.
\end{answerbox}
We show that $11$ cannot be reached by robot $B$. Suppose $11=3x+7y$ with whole numbers $x,y\ge0$.
\begin{itemize}
  \item If $y=0$, then $3x=11$, but $11$ is not a multiple of $3$.
  \item If $y=1$, then $3x=4$, which is also impossible.
  \item If $y\ge2$, then $7y\ge14>11$, too big.
\end{itemize}
So $B_{11}=0$ (this matches the table). If we picked any $N\le11$, then $n=11$ would have to work for both robots, but it does not.
So $N$ cannot be smaller than $12$, and we showed $12$ works. Hence the smallest $N$ is $12$.

\subsection{Showing $A_n+B_n\le(5/4)^n$ for all $n\ge10$}
Write $t=5/4$. The idea: we prove that each robot separately is at most \emph{half} of $t^n$. Then adding them gives at most $t^n$.
\[
A_n\le\tfrac12t^n\quad(n\ge4),\qquad B_n\le\tfrac12t^n\quad(n\ge6).
\]
Then for $n\ge10$ (in fact for $n\ge6$) we get $A_n+B_n\le\tfrac12t^n+\tfrac12t^n=t^n$.

\textbf{Step 1: Two small facts.}
\[
t^{-2}+t^{-5}=\tfrac{16}{25}+\tfrac{1024}{3125}=\tfrac{2000+1024}{3125}=\tfrac{3024}{3125}<1,
\]
\[
t^{-3}+t^{-7}=\tfrac{64}{125}+\tfrac{16384}{78125}=\tfrac{40000+16384}{78125}=\tfrac{56384}{78125}<1.
\]

\textbf{Step 2: The bound for $A_n$, by strong induction.}
$A_n$ looks back $2$ and $5$ steps, so we check $5$ starting values $n=4,5,6,7,8$.
The bound $A_n\le\frac12t^n$ is the same as $2A_n4^n\le5^n$ (multiply both sides by $2\cdot4^n$), which uses only whole numbers:
\begin{center}\small
\begin{tabular}{@{}c|ccccc@{}}
$n$ & 4 & 5 & 6 & 7 & 8\\\hline
$A_n$ & 1 & 1 & 1 & 2 & 1\\
$2A_n4^n$ & $512$ & $2048$ & $8192$ & $65536$ & $131072$\\
$5^n$ & $625$ & $3125$ & $15625$ & $78125$ & $390625$\\
True? & \cmark&\cmark&\cmark&\cmark&\cmark
\end{tabular}
\end{center}
Now let $n\ge9$ and assume the bound for all values from $4$ up to $n-1$. Since $n-2\ge7$ and $n-5\ge4$, both $A_{n-2}$ and $A_{n-5}$ satisfy the bound, so
\[
A_n=A_{n-2}+A_{n-5}\le\tfrac12t^{n-2}+\tfrac12t^{n-5}=\tfrac12t^n\,(t^{-2}+t^{-5})<\tfrac12t^n .
\]
(the last step uses Step 1). So the bound holds for all $n\ge4$.

\textbf{Step 3: The bound for $B_n$, by strong induction.}
$B_n$ looks back $3$ and $7$ steps, so we check $7$ starting values $n=6,\dots,12$:
\begin{center}\small
\begin{tabular}{@{}c|ccccccc@{}}
$n$ & 6 & 7 & 8 & 9 & 10 & 11 & 12\\\hline
$B_n$ & 1 & 1 & 0 & 1 & 2 & 0 & 1\\
$2B_n4^n$ & $8192$ & $32768$ & $0$ & $524288$ & $4194304$ & $0$ & $33554432$\\
$5^n$ & $15625$ & $78125$ & $390625$ & $1953125$ & $9765625$ & $48828125$ & $244140625$\\
True? & \cmark&\cmark&\cmark&\cmark&\cmark&\cmark&\cmark
\end{tabular}
\end{center}
Let $n\ge13$ and assume the bound for all values from $6$ up to $n-1$. Since $n-3\ge10$ and $n-7\ge6$,
\[
B_n=B_{n-3}+B_{n-7}\le\tfrac12t^{n-3}+\tfrac12t^{n-7}=\tfrac12t^n\,(t^{-3}+t^{-7})<\tfrac12t^n .
\]
So the bound holds for all $n\ge6$.

\textbf{Step 4: Put them together.} For every $n\ge10$,
\[
A_n+B_n\le\tfrac12\Big(\tfrac54\Big)^n+\tfrac12\Big(\tfrac54\Big)^n=\Big(\tfrac54\Big)^n. \qed
\]

% \begin{notebox}
% Why we used the factor $\tfrac12$: robot $A$ grows almost as fast as $(5/4)^n$ (its growth rate is about $1.2365$ and $5/4=1.25$), so there is not much room.
% Giving each robot exactly half of $(5/4)^n$ and checking the starting values carefully makes the proof work.
% \end{notebox}

\end{document}